\section{From tree excursions to the measured quantum sphere}
\label{tree:sec:sphere}

Throughout this section, $\gamma=\sqrt2$.

The local tree passage law has now been identified up to a deterministic
constant. We transfer this identification to a finite excursion, control
the two ends of its contour, and recover all unrestricted distances.
The contour parametrization will identify the limiting probability
measure at the same time.

Write $N=2n$ and let $Z^{\mathrm{disc}}$ be the walk whose independent
increments choose each of the four axial unit steps with probability
$1/4$.  Put $\mathcal Q=\Z_{\geq0}^2$ and
\[
 q_m(x,y)=\Prob_x[Z^{\mathrm{disc}}_m=y,\ Z^{\mathrm{disc}}_j\in\mathcal Q
                    \text{ for }0\leq j\leq m].
\]
Thus $q_m$ includes the killing condition.  We denote the law conditioned on
$Z^{\mathrm{disc}}_N=0$ and survival until $N$ by $\Prob_N^{\mathrm{exc}}$.
Theorem~\ref{tree:thm:background_mullin} identifies this law with the uniform pair
$(M_n,T_n)$ defined in Section~\ref{tree:sec:introduction}.  In particular,
every graph distance below uses all primal
edges of that pair.  The internal interval distances used as upper bounds
retain the tree edges and the matched non-tree edges witnessed in that
interval; they do not change the definition of the ambient distance.

We use the preliminary deterministic units
\begin{equation}\label{tree:eq:sphere_units}
 B_N=A_{\lfloor\log_2N\rfloor}.
\end{equation}
Theorem~\ref{tree:thm:bilateral_limit} says that $B_N\to\infty$, that every
sequence $N\to\infty$ has a further joint extraction of the bilateral
contours and passage data, and that on this extraction the local rule is
$\Lambda D_h$, for a deterministic $\Lambda$ in a fixed compact subset of
$(0,\infty)$.  Brownian time is normalized to have total duration one;
the deterministic factor from $N/2^{\lfloor\log_2N\rfloor}$ is included
in $\Lambda$.  We also use its uniform interval bound: if $I$ is an
integer interval with $1\leq |I|\leq N$, and $\mathcal D(I)$ is its
internal contour-vertex diameter, then
\begin{equation}\label{tree:eq:sphere_internal_bound}
 \big\|B_N^{-1}\mathcal D(I)\big\|_{L^k(\Prob)}
       \leq C\left(\frac{|I|}{N}\right)^s.
\end{equation}
A bound with $|I|+1$ is equivalent here.  These statements hold also on
any fixed finite enlargement of the time window.  No limit of adjacent
normalizing ratios will be assumed.

\subsection{Conditioning and the additional passage data}

\begin{lemma}[Excursion densities]\label{tree:lem:excursion_density}
There are constants $c,C>0$ such that, for all $m\geq0$ and $x\in\mathcal Q$,
\begin{equation}\label{tree:eq:sphere_entrance_bound}
 q_m(0,x)\leq C(x_1+1)(x_2+1)(m+1)^{-3}
                 \exp\left(-\frac{c|x|^2}{m+1}\right).
\end{equation}
For positive even $N$,
\begin{equation}\label{tree:eq:sphere_return}
 q_N(0,0)\asymp (N+1)^{-3},\qquad
       \lim_{N\to\infty,\,N\text{ even}}N^3q_N(0,0)=\frac{32}{\pi}.
\end{equation}
For $1\leq l\leq N/4$, the density of the increment path on $[l,2l]$
under $\Prob_N^{\mathrm{exc}}$, relative to the law of $l$ free increments,
is
\begin{equation}\label{tree:eq:sphere_endpoint_density}
 f_{N,l}(\omega)=\frac{1}{q_N(0,0)}\sum_{x\in\mathcal Q}q_l(0,x)
  \mathbf1_{\{x+\omega_j\in\mathcal Q,\ 0\leq j\leq l\}}
                      q_{N-2l}(x+\Delta,0),\qquad \Delta=\omega_l.
\end{equation}
It satisfies
\begin{equation}\label{tree:eq:sphere_polynomial_density}
 f_{N,l}(\omega)\leq C\prod_{i=1}^2
       \left(1+\frac{|\Delta_i|}{\sqrt l}\right)
       \leq C\left(1+\frac{|\Delta|}{\sqrt l}\right)^2,
 \qquad \sup_{N,l}\E[f_{N,l}^2]<\infty.
\end{equation}

Fix $0<a<b<1$.  The densities for the increment block
$[\lfloor aN\rfloor,\lfloor bN\rfloor]$ converge, under the free
invariance principle, to a continuous function of the limiting block's
coordinate minima and net increment.  They are uniformly integrable,
indeed uniformly bounded in $L^2$.  Consequently, if arbitrary additional
block-measurable passage data converge jointly with the free rescaled
increments along a subsequence, their excursion-conditioned joint laws
converge along that subsequence by the same change of density.  Conditional
on the limiting block increments, this change preserves their conditional
kernel.
\end{lemma}

\begin{proof}
Let $p_r(h)$ be the simple symmetric one-dimensional walk transition
probability from $0$ to $h$, and let $q_r^{(1)}$ denote its kernel killed
on entering $-1$.  Reflection gives, on the admissible parity lattice,
\begin{equation}\label{tree:eq:sphere_reflection_1d}
 q_r^{(1)}(0,h)=p_r(h)-p_r(h+2)
       =\frac{2(h+1)}{r+h+2}p_r(h),\qquad h\geq0.
\end{equation}
The formula is interpreted directly at $r=0$ and at the ends of the walk's
support.  The binomial formula for $p_r$, or Stirling's inequalities, gives
\[
 q_r^{(1)}(0,h)\leq C(h+1)(r+1)^{-3/2}
                    e^{-ch^2/(r+1)}.
\]
Given the number $J$ of first-coordinate moves in $m$ steps, the two
coordinate walks are independent and $J$ has the binomial law
$\mathrm{Bin}(m,1/2)$.  Thus
\begin{equation}\label{tree:eq:sphere_schedule_kernel}
 q_m(0,x)=2^{-m}\sum_{j=0}^m\binom mj
                  q_j^{(1)}(0,x_1)q_{m-j}^{(1)}(0,x_2).
\end{equation}
For $m/4\leq j\leq3m/4$, the preceding bound proves
\eqref{tree:eq:sphere_entrance_bound} term by term.  The complementary binomial
probability is at most $Ce^{-cm}$.  On the support of $q_m(0,x)$,
$|x|\leq m$; after decreasing $c$, this error is at most
$C(m+1)^{-3}e^{-c|x|^2/(m+1)}$ and is absorbed in the same bound.
The kernel is symmetric, so the bound also applies to $q_m(x,0)$.

Set $e_r=q_r^{(1)}(0,0)$.  The exact formula is
$e_{2j}=\binom{2j}{j}/(4^j(j+1))$, with $e_0=1$ and $e_{2j+1}=0$.
Hence $e_r\sim 2^{3/2}\pi^{-1/2}r^{-3/2}$ through even $r$.
Apply this to \eqref{tree:eq:sphere_schedule_kernel} with $x=0$ and $m=N$.
On $|J-N/2|\leq N^{3/4}$, the product of the two nonzero Catalan
probabilities is $(64/\pi+o(1))N^{-3}$, uniformly.  Exactly half the
binomial mass is on even $J$; the mass outside this window is exponentially
small compared with $N^{-3}$.  This proves the limit in
\eqref{tree:eq:sphere_return}; increasing the constants covers all positive
even $N$.

The Markov property at times $l$ and $2l$ proves
\eqref{tree:eq:sphere_endpoint_density}.  Since $N-2l\geq N/2$, use
\eqref{tree:eq:sphere_entrance_bound} twice and
\eqref{tree:eq:sphere_return}.  Drop the survival indicator and the exit
Gaussian.  Terms with $x+\Delta\notin\mathcal Q$ have zero weight, and
for the other terms $x_i+\Delta_i+1\leq x_i+|\Delta_i|+1$.  We obtain
\[
 f_{N,l}(\omega)\leq C l^{-3}\prod_{i=1}^2
  \sum_{h\geq0}(h+1)(h+|\Delta_i|+1)e^{-ch^2/l}.
\]
For $u\geq0$ and $l\geq1$, comparison with Gaussian integrals gives
\[
 \sum_{h\geq0}(h+1)(h+u+1)e^{-ch^2/l}
       \leq C\big(l^{3/2}+ul\big).
\]
This proves the pointwise bound in \eqref{tree:eq:sphere_polynomial_density}.
Free bounded increments have uniformly bounded moments of every fixed
order after division by $\sqrt l$: expanding an even moment, every
nonzero term pairs each independent centered increment with another
occurrence of that increment, giving a bound $C_r l^r$ for the $2r$th
moment.  In particular the fourth moment proves the asserted $L^2$ bound.
The same argument gives an $L^r$ bound for every fixed finite $r$.

Here are the local limits and the precise change-of-measure argument.
They also fix all parity factors.  Set $c_*=32/\pi$ and
\[
 K_t(u)=c_* t^{-3}u_1u_2e^{-|u|^2/t},\qquad t>0,\quad u\in\R_{\geq0}^2.
\]
Extend $K_t$ by zero outside the quadrant, and use the same zero
convention for killed kernels evaluated outside their state space.
Uniformly when $t$ stays in a compact subset of $(0,\infty)$ and $x/\sqrt N$
stays in a fixed compact subset of the quadrant, Stirling's formula in
\eqref{tree:eq:sphere_reflection_1d} and binomial concentration in
\eqref{tree:eq:sphere_schedule_kernel} imply
\begin{equation}\label{tree:eq:sphere_entrance_lclt}
 N^2q_{\lfloor tN\rfloor}(0,x)-K_t(x/\sqrt N)\longrightarrow0
\end{equation}
on the admissible lattice $x_1+x_2\equiv\lfloor tN\rfloor\pmod2$.
To check its coefficient, the one-dimensional local limit on its parity
lattice is $p_j(h)\sim\sqrt{2/(\pi j)}e^{-h^2/(2j)}$.
The product in \eqref{tree:eq:sphere_schedule_kernel}, with $j\sim tN/2$,
therefore has coefficient $64/\pi$; its admissible binomial schedules
have asymptotic mass $1/2$.  Boundary points cause no problem for uniform
absolute convergence: if one rescaled coordinate tends to zero, both
sides of \eqref{tree:eq:sphere_entrance_lclt} tend to zero by
\eqref{tree:eq:sphere_entrance_bound}.  This proves the stated uniform version.

Let $a_N=\lfloor aN\rfloor$, $b_N=\lfloor bN\rfloor$, and linearly
interpolate a free block after dividing heights by $\sqrt N$ and times
by $N$.  For its limiting path $w$ of duration $b-a$, write
$d=w(b-a)$ and $m_i=\min_t w_i(t)$.  The limiting density is
\begin{equation}\label{tree:eq:sphere_interior_density}
 f_{a,b}(w)=\frac1{2c_*}\int_{\R_{\geq0}^2}
 K_a(u)K_{1-b}(u+d)
       \mathbf1_{\{u_i+m_i\geq0,\ i=1,2\}}\,du.
\end{equation}
The discrete formula is \eqref{tree:eq:sphere_endpoint_density} with entrance
length $a_N$, block length $b_N-a_N$, and exit length $N-b_N$.
Each summation lattice has asymptotic density $N/2$ after rescaling;
\eqref{tree:eq:sphere_return} and \eqref{tree:eq:sphere_entrance_lclt} give
\eqref{tree:eq:sphere_interior_density}.  More explicitly, if rescaled paths
$w_N$ converge uniformly to $w$, their minima and increments converge.
On bounded $u$-sets the Riemann sums converge, since the boundaries
$u_i=-m_i$ have zero Lebesgue measure.  On their complements the Gaussian
bound \eqref{tree:eq:sphere_entrance_bound}, with bounded $d$, gives an
integrable common majorant.  It follows that $f_N(w_N)\to f_{a,b}(w)$.
The same argument proves continuity of $f_{a,b}$ in the uniform path
topology.  Dropping the exit Gaussian as above also gives
\[
 f_N(\omega)\leq C_{a,b}\left(1+\frac{|\Delta|}{\sqrt N}\right)^2.
\]
Thus the densities are uniformly bounded in $L^2$ under the free law.

For completeness, suppose that $(W_N,H_N)$ denotes the rescaled free
block together with any additional block data, taking values in the
Polish data spaces of Proposition~\ref{tree:prop:certificate_limit}, and that
$(W_N,H_N)\Rightarrow(W,H)$.  The preceding sequential convergence and
uniform integrability imply, for every bounded continuous $\Phi$,
\begin{equation}\label{tree:eq:sphere_weighted_joint_limit}
 \E\big[f_N(W_N)\Phi(W_N,H_N)\big]
       \longrightarrow \E\big[f_{a,b}(W)\Phi(W,H)\big].
\end{equation}
Indeed, realize the unweighted convergence almost surely, truncate the
densities at a fixed height, and then remove the truncation using their
uniform $L^2$ bound.  Taking $\Phi=1$ shows that the limiting density
integrates to one.  If the unweighted limit disintegrates as
$\mu(dw)\mathsf K(w,dz)$, its weighted law is exactly
$f_{a,b}(w)\mu(dw)\mathsf K(w,dz)$.  This proves the assertion about
additional data and their conditional kernel, not just convergence of
the contour marginal.
\end{proof}

The corresponding full-contour invariance principle has a useful direct
interpretation.  Given a schedule with $J$ first-coordinate moves,
excursion conditioning makes the two coordinate strings independent
nonnegative bridges of lengths $J$ and $N-J$.  The conditional mass of $J$
is
\begin{equation}\label{tree:eq:sphere_schedule_excursion}
 \frac{2^{-N}\binom NJ e_Je_{N-J}}{q_N(0,0)}.
\end{equation}
It is concentrated at $J/N=1/2$.  For each $\epsilon>0$, the probability
of $|J-N/2|>\epsilon N$ is at most
$C N^3 e^{-c_\epsilon N}$, by the ordinary binomial bound and
\eqref{tree:eq:sphere_return}.  Conditional on $J$, all schedules are uniform
and independent of the two coordinate bridges.  Their partial counts
satisfy \mbox{$\sup_{r\leq N}|J_r-rJ/N|/N\to0$} in probability.  One way to
see the last assertion is to use
\[
 \Prob[J_r=u\mid J]=\binom ru\binom{N-r}{J-u}\big/\binom NJ
\]
and Stirling's inequalities: for $J/N$ in a compact subset of $(0,1)$,
the mass with $|u-rJ/N|>\epsilon N$ is at most
$C N^C e^{-c_\epsilon N}$, uniformly in $r$; summing over $r$ suffices.
Consequently the schedule time changes converge uniformly to half speed.
The reflection formula gives the nonnegative-bridge finite-dimensional
limits as the Brownian excursion transition densities.  It also gives
tightness: on a fixed interior time interval use the bounded $L^2$
density with respect to free increments; close to each endpoint use the
one-dimensional version of \eqref{tree:eq:sphere_polynomial_density} and
sum dyadic free-walk maximum bounds.  The latter sum is bounded by
$C\sum_{j\geq0}\sqrt{\delta2^{-j}}=C'\sqrt\delta$.
This proves the excursion invariance principle, including the endpoints.
Thus $N^{-1/2}Z^{\mathrm{disc}}_{\lfloor Nt\rfloor}$ converges to two
independent Brownian excursions with the variance convention inherited
from the axial walk.

By Theorem~\ref{tree:thm:background_peanosphere}, equivalently
\cite[Theorem~1.1]{MillerSheffieldSphere}, this pair encodes the area-one quantum
sphere, with its area-parametrized space-filling $\SLE_8$.  This is the
ordinary sphere of \cite[Definition~4.21(ii)]{DMS}.  Its distinguished
points can be forgotten, by \cite[Proposition~A.13]{DMS}; we write $q$
for the common point $\eta(0)=\eta(1)$ used by the encoding.  The
strict positivity of both excursions on $(0,1)$ gives
\begin{equation}\label{tree:eq:sphere_single_endpoint}
                      \eta^{-1}(\{q\})=\{0,1\}.
\end{equation}
All deterministic Brownian boundary-length conventions here are the same
ones as in the bilateral construction and
Theorem~\ref{tree:thm:background_peanosphere}.

\subsection{The metric on interior sphere charts}

\begin{lemma}[Interior sphere metric]\label{tree:lem:sphere_interior_metric}
In the joint extraction with distance units $B_N$ and deterministic
multiplier $\Lambda$ selected above, the limiting finite-excursion
local passage rule satisfies
\[
 F_U(x,y)=\Lambda D_h(x,y;U)
\]
for every buffered chart $U$ compactly contained in $S\setminus\{q\}$
and every $x,y$ joined by a path compactly in $U$. This identification
retains the upper realization and lower subpath-extraction conclusions
of Theorem~\ref{tree:thm:bilateral_limit}.
\end{lemma}
\begin{proof}
Fix rational $0<\delta<1/2$, and write
$a=\delta$, $b=1-\delta$. Retain the complete
collection of internal interval-chart passage data constructed solely
from the increments on $[a_N,b_N]$.  It includes every certificate whose
range has positive clearance from the boundary of the interval surface,
by Lemma~\ref{tree:lem:chart_reconstruction}.  Apply
\eqref{tree:eq:sphere_weighted_joint_limit} to this collection, first in
finitely many compact charts and then by the countable extraction of
Proposition~\ref{tree:prop:certificate_limit}.  In the bilateral reference
law its compact-interior local passage rule equals $\Lambda D$, by
Theorem~\ref{tree:thm:bilateral_limit}; the equality event has probability
one and is therefore preserved by the limiting density.

The restricted increments determine an abstract decorated quantum
surface, including its quantum boundary-length units. For the cone
this is \cite[Lemma~4.9]{GHSSurvey}. The sphere statement needed here
is supplied separately by the final paragraph of the proof of
\cite[Theorem~1.1, the case $\gamma\leq\sqrt2$, p.~49]{MillerSheffieldSphere}:
conditionally on the contour values at $\delta$ and $1-\delta$, the
quantum surface complementary to
$\eta([0,\delta])\cup\eta([1-\delta,1])$ has the corresponding cone
conditional law and is determined, together with its curve, by the
restricted contour. This is the symmetric interval statement recorded
in Theorem~\ref{tree:thm:background_peanosphere}. Its quantum units agree
with those in the density calculation. Metric locality and conformal
covariance therefore identify its internal $D$ with the internal metric
on its image in the area-one sphere. The transported equality is an
equality with this sphere's $D_h$ in the fixed normalization.

This equality is an event of the restricted contour and its passage
data, since those contour increments determine the abstract interval
surface. Adjoining the outside contour therefore preserves the same
probability-one equality.

Finally, if $K\subset S\setminus\{q\}$ is compact, continuity of $\eta$
and \eqref{tree:eq:sphere_single_endpoint} imply that
$\eta^{-1}(K)$ is contained in $(\delta,1-\delta)$ for some rational
$0<\delta<1/2$. Choose a smaller positive rational $\delta$ if necessary.
The image of the omitted
compact time intervals is disjoint from $K$, so $K$ lies in the interior
of the interval surface and has positive spatial clearance from its
boundary.  Its local certificates are therefore among the ones just
transferred. By a countable exhaustion, the local passage rule off $q$
equals $\Lambda D_h$ simultaneously in these charts, with the same
deterministic $\Lambda$.

\end{proof}

We also record how the transfer behaves when the complete outside
contour is retained. The following conditional-kernel statement allows
the inside passage observations to be used jointly with outside
contour observations.

\begin{lemma}[The kernel after adjoining the outside contour]
\label{tree:lem:sphere_exterior_kernel}
Fix $0<a<b<1$. In the setup of
Equation~\eqref{tree:eq:sphere_weighted_joint_limit}, suppose the free
joint limit of the block increments $W_N$ and block-measurable passage
data $H_N$ disintegrates as $\mu(dw)\mathsf K(w,dz)$.
Write $W,H$ for the corresponding interior increment path and passage
datum in their joint excursion-conditioned limit with the complete
contour. The conditional law of $H$ given that contour is
$\mathsf K(W,\cdot)$.
\end{lemma}
\begin{proof}
The assertion follows from the Markov factorization: given
the entrance height $x$ and the inside increments, the outside consists
of a killed entrance bridge and a killed exit bridge, independent of
the inside data.  Their weights are the two factors in
\eqref{tree:eq:sphere_endpoint_density}.  On compact positive-height ranges,
the reflection formula for $q_m(x,y)$ (the four coordinate reflections
of the free two-dimensional kernel) gives uniform bridge
finite-dimensional limits.  For tightness at the zero endpoint, the
one-dimensional reflection formula gives
$q_r^{(1)}(u,v)\leq C(u+1)(v+1)(r+1)^{-3/2}$: telescope step-two
differences of $p_r$ in the smaller of $u,v$, using symmetry of the
kernel and the binomial Gaussian bound.  Averaging over schedules as in
\eqref{tree:eq:sphere_schedule_kernel} consequently gives
$q_m(y,z)\leq C_R(y_1+1)(y_2+1)N^{-2}$ when $m\asymp N$ and
$|z|\leq R\sqrt N$.  On a fixed positive-height compact set,
$q_{a_N}(0,z)\geq c_RN^{-2}$, so the proof of
\eqref{tree:eq:sphere_polynomial_density} applies to the small dyadic blocks
of this entrance bridge.  At its other endpoint, reverse the bridge:
the increment density on an initial short block is bounded by
$C_R(1+|\Delta|/\sqrt N)^2$, using
\eqref{tree:eq:sphere_entrance_bound}.  Free-walk maximum bounds therefore
give tightness there too.  This proves the uniform bridge convergence
used here.

Both entrance and exit heights can be restricted to positive compact
sets before this calculation.  Indeed their marginal law at time $j$ is
\[
 \Prob_N^{\mathrm{exc}}[Z_j^{\mathrm{disc}}=x]
       =\frac{q_j(0,x)q_{N-j}(x,0)}{q_N(0,0)}.
\]
Uniformly for $j/N$ in a compact subinterval of $(0,1)$,
\eqref{tree:eq:sphere_entrance_bound} bounds this by
$CN^{-3}(x_1+1)^2\allowbreak\mbox{$(x_2+1)^2$}e^{-c|x|^2/N}$.
Summing shows that the probability of either coordinate below
$\epsilon\sqrt N$ has limsup at most $C\epsilon^3$, whereas the
probability of $|x|>R\sqrt N$ tends to zero uniformly as $R\to\infty$.
Apply these bounds at both $a_N$ and $b_N$.

Retain the rescaled entrance height $u=x/\sqrt N$ jointly in the Riemann
sum proving \eqref{tree:eq:sphere_weighted_joint_limit}.  On the preceding
compact sets, the two bridge kernels converge uniformly and are
independent of the additional inside data given $(u,w)$.  The limiting
joint law consequently has the form
\[
 \mu(dw)\mathsf K(w,dz)\,
 \frac{K_a(u)K_{1-b}(u+d)}{2c_*}
 \mathbf1_{\{u_i+m_i\geq0,\ i=1,2\}}\,du\,
 \mathsf B_a(0,u;d\omega^-)\mathsf B_{1-b}(u+d,0;d\omega^+),
\]
where $\mathsf B$ are the limiting killed bridge laws.  The marginal
tail bounds just proved remove the truncations.  Conditioning on the
complete contour $(\omega^-,w,\omega^+)$ therefore leaves the inside
kernel $\mathsf K(w,dz)$ unchanged, exactly as in
Lemma~\ref{tree:lem:finite_kernel_transfer}.

\end{proof}

\subsection{Endpoint diameters}

The density estimate controls the metric information which an interior
change of measure cannot see.  Let $v_i$ be the primal vertex at contour
time $i$, with $v_0=v_N$, and set
\[
 T_N(\delta)=\{v_i:0\leq i\leq\lfloor\delta N\rfloor
                   \text{ or }N-\lfloor\delta N\rfloor\leq i\leq N\}.
\]
For $1\leq l\leq N/4$, Cauchy--Schwarz,
\eqref{tree:eq:sphere_polynomial_density}, and
\eqref{tree:eq:sphere_internal_bound} give
\begin{equation}\label{tree:eq:sphere_endpoint_block}
 \E_N^{\mathrm{exc}}\left[B_N^{-1}\mathcal D([l,2l])\right]
 \leq \|f_{N,l}\|_2
         \left\|B_N^{-1}\mathcal D([0,l])\right\|_2
 \leq C(l/N)^s.
\end{equation}
The internal distances in this display are genuine functions of the block
increments, so the density calculation applies to them exactly.
If only endpoint norms are used in
Theorem~\ref{tree:thm:bilateral_limit}, the diameter bound follows by balanced
binary subdivision: at depth $r$, the maximum of at most $2^r$ endpoint
costs has $L^k$ norm at most
$C(l/N)^s2^{-r(s-1/k)}$ in $B_N$ units; summing these bounds is finite
since $s>1/k$.  Rounded subdivision lengths change the constant only.

For $0<\delta\leq1/8$ and $\delta N\geq1$, choose $J$ with
$2^J\leq\delta N<2^{J+1}$.  The chain
\[
 [0,1],\ [1,2],\ [2,4],\ldots,\ [2^J,2^{J+1}]
\]
covers $[0,\delta N]$ and its consecutive intervals share endpoints.
Every block except $[0,1]$ in this chain satisfies the hypothesis of
\eqref{tree:eq:sphere_endpoint_block}.  Its first interval has graph
diameter at most one.  Connect points through the common endpoints
using the interval paths and sum the resulting diameter bounds.
Time reversal gives the same bound at $N$, and the two chains meet
at $v_0=v_N$.  If $\delta N<1$, the endpoint set consists only of this
vertex.  We have proved
\begin{equation}\label{tree:eq:sphere_no_spikes}
 \E_N^{\mathrm{exc}}\left[B_N^{-1}\diam_{d_n}T_N(\delta)\right]
       \leq C\delta^s+2B_N^{-1},\qquad 0<\delta\leq1/8.
\end{equation}
In particular,
\begin{equation}\label{tree:eq:sphere_no_spikes_probability}
 \lim_{\delta\downarrow0}\limsup_{N\to\infty}
 \Prob_N^{\mathrm{exc}}\left[
       B_N^{-1}\diam_{d_n}T_N(\delta)>\epsilon\right]=0
       \quad\text{for every }\epsilon>0.
\end{equation}
This is an estimate for the actual ambient graph metric; the internal
metrics appeared only to construct admissible paths.

On a closed interior time interval, the uniform $L^2$ density in
Lemma~\ref{tree:lem:excursion_density} transfers the free time-modulus
estimate of Lemma~\ref{tree:lem:moment_compactness}.  Explicitly, if an event
has free probability $r$, its conditioned probability is at most
$C_{a,b}\sqrt r$.  Together with
\eqref{tree:eq:sphere_no_spikes_probability}, this gives tightness of the
cyclic time modulus of the rescaled ambient distances on $[0,1]$.
Indeed first cover $[\delta,1-\delta]$ by finitely many interior blocks,
and then join its ends to the two endpoint chains just constructed.
The same argument bounds the full diameter in probability.

\subsection{Identification of unrestricted distances}

\begin{lemma}[The measure counted by the tree tour]
\label{tree:lem:tour_measure}
Let $v_i$ be the primal vertex represented by contour time $i$, with
$v_N=v_0$. Then, for every $v\in V(M_n)$,
\begin{equation}\label{tree:eq:tour_degree}
 \#\{0\leq i<N:v_i=v\}=\deg_{M_n}(v).
\end{equation}
Consequently, if
\[
 \mu_n=\frac1{2n}\sum_{v\in V(M_n)}\deg_{M_n}(v)\,\delta_v,
 \qquad f_n(t)=v_{\lfloor Nt\rfloor},\quad 0\leq t\leq1,
\]
then $(f_n)_*\mathrm{Leb}_{[0,1]}=\mu_n$. Conditional on the map,
independent uniform times give independent samples from $\mu_n$.
\end{lemma}

\begin{proof}
Associate time slot $i\in\{0,\ldots,N-1\}$ with the step departing
at $i$. The two horizontal traversals of a tree edge depart once
from each endpoint, so these slots count its two incidences. Every
vertical step belongs to a unique matched pair with terminal times
$u<v$. By \eqref{tree:eq:discrete_R_match}, this pair adds a non-tree
edge between $v_{u-1}$ and $v_{v-1}$; its two departure slots count
exactly those two incidences. This also counts a loop twice. Every
slot belongs to exactly one such pair, proving
\eqref{tree:eq:tour_degree}. Each slot has time mass $1/N=1/(2n)$,
which proves the pushforward identity and the sampling assertion.
\end{proof}

\begin{proposition}[Sphere identification]\label{tree:prop:sphere_identification}
For every sequence of positive integers $n\to\infty$, there is a further
subsequence and a deterministic $\Lambda\in(0,\infty)$ such that
\begin{equation}\label{tree:eq:sphere_preliminary_limit}
        (V(M_n),B_{2n}^{-1}d_n,\mu_n)
              \ \Rightarrow\ (S,\Lambda D_h,\mu_h)
\end{equation}
in the Gromov--Hausdorff--Prokhorov topology. The constants $\Lambda$ lie in a
fixed deterministic compact subset of $(0,\infty)$.
\end{proposition}

\begin{proof}
Choose a simultaneous bilateral extraction as in
Theorem~\ref{tree:thm:bilateral_limit}, with its deterministic multiplier
$\Lambda$. Lemma~\ref{tree:lem:sphere_interior_metric} gives convergent
sphere certificates on every compact subset of $S\setminus\{q\}$,
with local passage rule $\Lambda D_h$. The local
certificate assertions include realization with open margins, subpath
extraction, and uniformly negligible costs for joining nearby endpoints.

We give the compactness and path argument explicitly.  Put
\[
 \rho_N(t,u)=B_N^{-1}d_n(v_{\lfloor Nt\rfloor},v_{\lfloor Nu\rfloor}),
                    \qquad 0\leq t,u\leq1.
\]
Consecutive contour vertices are equal or joined by a tree edge, so
replacing $\rho_N$ by its bilinear interpolation changes it by at most
$2/B_N$.  The cyclic time modulus and diameter bound proved above imply
tightness of these interpolations in $C([0,1]^2)$.  Extract jointly with
the contours and all certificate charts, and use an almost sure
representation.  Then $\rho_N$ converges uniformly to a continuous
pseudometric $\rho$ on $[0,1]$, with $\rho(0,1)=0$.  The limiting
contour determines the sphere and its continuous surjection $\eta$.
In particular the endpoint time sets have limiting $\rho$-diameters
tending to zero.  We will show that
\begin{equation}\label{tree:eq:sphere_uniform_identification}
                       \rho(t,u)=\Lambda D_h(\eta(t),\eta(u))
                       \quad(0\leq t,u\leq1).
\end{equation}

The projection part of Proposition~\ref{tree:prop:certificate_limit} applies
to the finite excursion sewing as well: its proof uses uniform contour
convergence and the defining link relations.  Thus all representatives
of a single vertex have $D_h$-close images, uniformly over vertices, and
the two endpoints of every graph edge have $D_h$-distance tending to
zero, uniformly over edges.  To see the uniform assertion directly,
a contrary sequence has convergent contour times; the limiting tree
identification or matched second-coordinate chord identifies the two
images.  This also covers links involving time $0$ or $1$, which
project to $q$.  Choose any representative of each vertex and write
$\pi_N(v)$ for its image under $\eta$ at the rescaled time.  Let
$\varepsilon_N\to0$ bound both representative errors and edge-image
jumps.  These maps are only used to locate paths in the limiting
surface; all costs remain discrete all-edge costs.

We first record a uniform local crossing fact.  If $K\subset S\setminus\{q\}$
is compact, take finitely many pairs of open sets $V_i\Subset U_i$
covering $K$, all contained in interval-surface interiors.  Choose
$r>0$ so small that every $D_h$-ball of radius $3r$ about a point of $K$
is contained in one of the $U_i$ with positive clearance from its
boundary.  A discrete path beginning at such a point, stopped on first
reaching $D_h$-displacement $r$, stays in that compact chart after an
$\varepsilon_N$ enlargement.  For all sufficiently large $N$, it has
scaled cost at least $\Lambda r/2$.  Otherwise a subsequence of these
paths has bounded cost and compact confinement; subpath extraction
gives a certificate with endpoints at distance $r$ and cost at most
$\Lambda r/2$, contradicting the local lower comparison.  Compactness
of $K$ and the finite cover make this assertion uniform in its starting
point and in the path.  Paths of larger cost already satisfy it.
It follows that a path of scaled cost at most $L$ can be divided into
at most $2L/(\Lambda r)+2$ such steps, with a final shorter step, as
long as it stays in $K$.

For the upper bound in \eqref{tree:eq:sphere_uniform_identification}, fix
$t,u$ and a $D_h$-geodesic $P$ from $x=\eta(t)$ to $y=\eta(u)$.
A compact subarc disjoint from $q$ has a finite subdivision into pieces
lying with open margins in the sets $U_i$ of a finite local cover.
For each piece the local upper comparison and realization supply
approximating discrete paths with cost at most $\Lambda$ times that
piece's $D_h$-length, plus an arbitrarily prescribed positive error.
Uniform joining connects their endpoints, also with vanishing cost.
Consequently a geodesic whose range avoids $q$ gives
$\rho(t,u)\leq\Lambda D_h(x,y)$.

Here is the endpoint modification when $P$ meets $q$, also valid if one
or both endpoints are $q$.  Given $\delta>0$, compactness and
\eqref{tree:eq:sphere_single_endpoint} give a radius $r_\delta>0$ such that
\begin{equation}\label{tree:eq:sphere_endpoint_ball}
 \eta^{-1}(\overline B_{D_h}(q,2r_\delta))
                         \subset[0,\delta)\cup(1-\delta,1].
\end{equation}
If $P$ enters $B_{D_h}(q,r_\delta)$, retain its portion before the first
entry and its portion after the last exit.  Both are compactly away
from $q$ except for their retained endpoints, which are still at
positive distance from $q$ unless the whole portion is empty.  Apply
the finite local subdivision to these portions.  The two cut endpoints
have contour representatives in the endpoint time set of
\eqref{tree:eq:sphere_endpoint_ball}.  Their approximating vertices can
be connected at cost at most the scaled diameter of $T_N(\delta)$,
with a vanishing joining error; the same applies if an original endpoint
lies in this set.  The sum of the retained geodesic lengths is at most
$D_h(x,y)$.  Uniform convergence of $\rho_N$ therefore gives
\[
 \rho(t,u)\leq\Lambda D_h(x,y)+
       \diam_\rho([0,\delta]\cup[1-\delta,1]).
\]
The last term tends to zero by continuity of $\rho$ at $0$ and $1$.
This proves the upper bound in all cases.

For the lower bound, choose shortest discrete paths between the two
contour vertices.  Their scaled lengths converge to $\rho(t,u)$ and
are bounded.  Fix $r>0$.  If a path has vertices whose images lie in
$B_{D_h}(q,r)$, remove the portion between its first and last such
vertices, retaining its prefix and suffix.  The removed portion has
nonnegative cost.  Its two cut endpoints have $D_h$-distance at most
$2r$.  Apart from these endpoints, the retained pieces stay outside
$B_{D_h}(q,r)$; since edge-image jumps are at most $\varepsilon_N$,
they are contained in $S\setminus B_{D_h}(q,r/2)$ for large $N$.
If there is no visit, retain the entire path.  Cases with an endpoint
inside the ball simply have an empty retained piece.

Cover $S\setminus B_{D_h}(q,r/2)$ by the finite buffered local charts
above and choose a fixed step size $r'>0$ for that cover.  Divide each
retained piece at successive first $D_h$-displacements $r'$.
The uniform crossing fact bounds the number of pieces independently
of $N$.  Pass to a further extraction recording their endpoints,
costs, and marked subpaths.  Each piece has a limiting certificate
compactly in one chart.  Its limiting cost is at least $\Lambda$
times its endpoints' $D_h$-distance.  Add these inequalities, including
the final shorter pieces.  The triangle inequality in $D_h$ then gives
\[
 \rho(t,u)\geq \Lambda\big(D_h(x,y)-2r\big).
\]
No cost is assigned to the removed segment in this bound.  Letting
$r\downarrow0$ proves the required lower inequality.  In particular,
it rules out a short global path escaping a local chart and returning
at a different point; such a path has paid for its successive chart
crossings.  We have proved \eqref{tree:eq:sphere_uniform_identification}.

The construction applies first on a countable dense set of contour
pairs.  Continuity of $\rho$, $D_h$, and $\eta$ extends it to every
pair.  As $\rho_N\to\rho$ uniformly, the correspondence
\[
 \mathcal R_N=\{(v_{\lfloor Nt\rfloor},\eta(t)):0\leq t\leq1\}
\]
has distortion tending to zero when the discrete metric is $B_N^{-1}d_n$
and the sphere metric is $\Lambda D_h$.  It covers both spaces: every
primal vertex has a contour visit and $\eta$ is surjective. Moreover,
Lemma~\ref{tree:lem:tour_measure} gives $(f_n)_*\mathrm{Leb}=\mu_n$,
and quantum-area parametrization gives $\eta_*\mathrm{Leb}=\mu_h$.
Lemma~\ref{measure:parameter} therefore bounds the GHP distance by
half the same distortion. This proves
\eqref{tree:eq:sphere_preliminary_limit}.  The bounds and deterministic
nature of $\Lambda$ come from Theorem~\ref{tree:thm:bilateral_limit}; a change
of density cannot change a fixed almost sure equality constant.
\end{proof}

\subsection{A deterministic normalization for every integer}

\begin{proof}[Proof of the spanning-tree case of Theorem~\ref{tree:thm:main}]
Let $X=\diam_{D_h}S$. The ordinary area-one quantum sphere is compact
and nontrivial, so $0<X<\infty$ almost surely. With the center $m$ from
Lemma~\ref{measure:center}, define, for every positive integer $n$,
\begin{equation}\label{tree:eq:sphere_final_normalization}
 X_n=\diam_{d_n}V(M_n),\qquad
 a_n=\exp\left\{m(\log X)-m\big(\log(X_n\vee1)\big)\right\}.
\end{equation}
These are positive deterministic constants; the truncation includes
one-vertex maps. Proposition~\ref{tree:prop:sphere_identification}
supplies every hypothesis of Lemma~\ref{measure:normalization}, with
$s_n=B_{2n}\to\infty$ and a deterministic positive subsequential
multiplier $\Lambda$. Hence $a_n\to0$ and
\[
 (V(M_n),a_nd_n,\mu_n)\Longrightarrow(S,D_h,\mu_h)
\]
in GHP topology along all positive integers. All distances use every
primal edge, and forgetting the root and spanning tree gives the
unmarked metric probability spaces in the theorem.
\end{proof}
